\documentclass[12pt, a4paper, twocolumn]{article}
\usepackage{fontspec}
\usepackage{xeCJK}
\usepackage{geometry}
\usepackage{titlesec}
\usepackage{fancyhdr}
\usepackage{parskip}
\usepackage{indentfirst}
\usepackage{booktabs}
\usepackage{cite}
\usepackage{amsmath}
\usepackage{amssymb}
\usepackage{xcolor}

\usepackage[
    unicode=true,
    colorlinks=true,
    linkcolor=blue!60!black,
    citecolor=blue!60!black,
    urlcolor=blue!60!black
]{hyperref}

% 設定頁面邊距
\geometry{
    a4paper,
    left=1cm,
    right=1cm,
    top=2cm,
    bottom=2cm
}

% 設定字型
\setmainfont[
    Path = ./fonts/Times_New_Roman/,
    BoldFont = *-bold,
    ItalicFont = *-italic,
    BoldItalicFont = *-bold_italic
]{times}

\setmonofont[
    Path = ./fonts/Roboto_Mono/static/,
    UprightFont = *-Regular,
    BoldFont = *-Bold,
    ItalicFont = *-Italic,
    BoldItalicFont = *-BoldItalic
]{RobotoMono}

\xeCJKsetup{AutoFakeBold=true, AutoFakeSlant=true}
\setCJKmainfont[
    Path = ./fonts/kaiu/,
]{kaiu}
\setCJKmonofont{./fonts/kaiu/kaiu.ttf}

% 頁面設定
\setlength{\headheight}{14.5pt}
\fancyhf{}
\rfoot{\thepage}

% 標題樣式
\titleformat{\section}{\large\bfseries}{\Roman{section}}{1em}{}
\titleformat{\subsection}{\normalsize\bfseries}{\thesubsection}{1em}{}

% 段落縮排
\setlength{\parindent}{2em}

\pagestyle{fancy}
\fancyhf{}
\rhead{June. 2026}
\lhead{\textbf{Speech and Language Processing}}
\rfoot{\thepage}


\newcommand{\crab}{\textsc{Crab}}

\begin{document}

\twocolumn[
\begin{@twocolumnfalse}
\begin{center}
    \LARGE \textbf{Speech and Language Processing Project Report}\\
    \large 簡子霖\\
    \normalsize 逢甲大學資訊工程所\\
    \vspace{0.4em}
    \normalsize \texttt{M1406785@o365.fcu.edu.tw}
\end{center}

\vspace{0.8em}
\end{@twocolumnfalse}
]

\section*{Abstract}

This study repurposes \crab{} (Contrastive Representation and Multimodal Aligned Bottleneck) from its original speech emotion recognition setting to audio deepfake detection.
The goal of audio deepfake detection is to determine whether a speech utterance is bona fide or spoofed by text-to-speech synthesis, voice conversion, replay attacks, or other speech generation techniques.
Unlike emotion recognition, which predicts affective categories, audio deepfake detection is a security-oriented binary countermeasure task;
therefore, the model must not only classify correctly but also produce robust spoofing scores under unseen attacks, unseen speakers, and cross-corpus conditions.

The original \crab{} paper builds a bimodal Cross-Modal Transformer with WavLM and RoBERTa,and introduces Multi Layer Contrastive Supervision (MLCS) to apply multi-positive contrastive learning at several intermediate layers, encouraging discriminative speech and text representations before the final classifier \cite{ueda2026crab}.
This report transfers the same representation-learning idea to audio deepfake detection: a self-supervised speech encoder extracts waveform representations, the prediction target is changed to bona fide/spoof classification, and evaluation follows ASVspoof-style metrics such as Equal Error Rate (EER) and detection cost functions.
The front part of this report focuses on the motivation, related work, backend detection system, and the role of \crab{} as a representation-learning basis for audio deepfake detection.

\noindent\textbf{Keywords:} audio deepfake detection, ASVspoof, self-supervised speech representation, contrastive learning, \crab

\vspace{1em}

\section{INTRODUCTION}
With the rapid advancement of neural text-to-speech (TTS) and voice conversion (VC) technologies, synthetic speech has become increasingly natural, speaker-similar, and emotionally expressive.
Although these technologies have beneficial applications in assistive communication, content creation, and voice interaction systems, they also reduce the cost of generating fake audio.
Attackers can use only a small amount of speaker data to synthesize speech that mimics specific individuals, enabling scams, forged voice evidence, and attacks against voice-based authentication systems.
Therefore, Audio Deepfake Detection (ADD) has become a critical research problem at the intersection of speech processing and information security.

In recent years, self-supervised learning (SSL) models have become important front-end components for speech-related tasks because they can learn rich representations from large amounts of unlabeled data.
Models such as wav2vec 2.0 \cite{baevski2020wav2vec2}, HuBERT \cite{hsu2021hubert}, and WavLM \cite{chen2022wavlm} have shown strong performance across various downstream tasks, including automatic speech recognition, speaker verification, emotion recognition, and spoofing detection.

In anti-spoofing research, recent end-to-end baselines directly learn discriminative cues from raw waveforms.
For example, the AASIST \cite{jung2021aasist} model uses an integrated spectro-temporal graph attention network to jointly model spoofing artifacts in both spectral and temporal regions, replacing the score-level ensemble of multiple subsystems with a single model.
However, using a large SSL encoder alone does not guarantee that the detector will have sufficient generalization ability.
The downstream classification head must still organize intermediate-layer information effectively and prevent the model from relying only on spoofing artifacts that are specific to the training data.

% crab model introduction
In this work, we use \crab{} as the backbone model.
Because audio deepfake detection does not mainly depend on semantic relevance, we replace the original text encoder with a prosody encoder.
The main contributions of this work are summarized as follows:
\begin{enumerate}
    \item We redefine the downstream task of the \crab{} model from Speech Emotion Recognition (SER) to Audio Deepfake Detection (ADD).
    \item We introduce a prosody encoder to provide additional prosodic features for improving deepfake detection.
    \item We conduct extensive experimental evaluation and analysis, and explain why MLCS and contrastive representation learning may help audio deepfake detection.

\end{enumerate}

This report is organized as follows.
Section 2 reviews related works on self-supervised representations, Previous ADD Method, multimodal ADD, and crab model. 
Section 3 describe the modification from crab to fetch ADD task, model architecture, and experience setup.
Section 4 present the experiment result and evaluation metrices.
Finally, Section 5 sum up the conclusion of the paper.

\section{RELATED WORK}
\subsection{Audio Deepfake Detection}
\subsection{Prosody feature in Audio Deepfake}

\subsection{\crab 模型}

\section{METHODOLOGY}
\subsection{系統架構}

\section{RESULT}

情緒辨識的評估輸出 WAR 與 UAR：
\begin{align}
    \mathrm{WAR}
    &= \frac{\sum_i \mathbb{I}(\hat{y}_i=y_i)}{N}, \\
    \mathrm{UAR}
    &= \frac{1}{K}\sum_{c=1}^{K}
    \frac{\sum_i \mathbb{I}(\hat{y}_i=c, y_i=c)}
    {\sum_i \mathbb{I}(y_i=c)}.
\end{align}
深偽語音偵測則以 softmax spoof score 計算 EER，並在產生 `evaluation\_scores.txt' 後呼叫 ASVspoof 評測模組計算 minDCF、EER、CLLR 與 Act-DCF。

\section{CONCLUSION}

\bibliographystyle{unsrt}
\bibliography{references}

\end{document}
